\subsection{Why the integration bound is conditional}
In Section~\ref{sec:numeval} we used a series near zero and two Gauss rules on the rest of
the interval. On all three subintervals we must distinguish the error of the integration
method from the error in the sampled function values. Cauchy bounds control the coefficients
that finite sampling combines or omits.

On the first interval $[0,r]$, with $r=0.03$, we write $k_6(x)=\sum k_jx^j$ and
$\mathcal V(x)=\sum v_jx^j$. The two vanishing coefficients of $k_6$ give the exact integral
$\sum_{j\ge2}k_jr^{j-1}/(j-1)$. The calculation extracts the Fourier coefficients of $d$ and
$\mathcal V$ separately. If a tilde denotes coefficients reconstructed from $N=32$ values on
$|x|=r$, then for $2\le j<N$ we have
\[
 \widetilde d_j+(1-j)\widetilde v_j-k_j
   =\sum_{\ell\ge1}\bigl(k_{j+\ell N}+\ell N v_{j+\ell N}\bigr)r^{\ell N}
\]
for exact samples. Indeed, sampling an analytic function combines the indices congruent to $j$
modulo $N$, and $k_m=d_m+(1-m)v_m$ for $m\ge1$. Each difference is integrated with weight
$r^{j-1}/(j-1)$, and the tail of the exact integral begins at $j=N$. Cauchy's formula on the
outer circle bounds both types of coefficients. The intermediate circle of radius $0.055$
uses 33 independent samples (64 including their conjugates) to improve the bounds. The largest
radius used in this estimate is $0.095$.

On the other two intervals we want to know which coefficients can contribute to the Gauss
quadrature error. On a real panel $[a,b]$, let $h=(b-a)/2$, $y=(a+b)/2$, and $f=k_6/x^2$. Write
\[
 g(w)=f\!\left(y+\tfrac h2(w+w^{-1})\right)
      =\sum_{j\in\mathbb Z}\beta_jw^j .
\]
The substitution $w\mapsto w^{-1}$ leaves the argument of $f$ unchanged, so $\beta_{-j}=\beta_j$.
By analyticity in the Bernstein ellipse of parameter $\rho_2$ (see~\cite{trefethen}), Cauchy's
formula gives $|\beta_j|\le\overline M\rho_2^{-j}$ for $j\ge0$, where $\overline M$ bounds the
mean modulus on the contour. Let $E_s$ be the integral minus the $s$-point Gauss rule on $[-1,1]$.
For the Chebyshev polynomial $T_j$ of degree $j$, the symmetry and exactness of the rule give
$E_s[T_j]=0$ if $j$ is odd or $j<2s$. For even $j\ge2s$, positivity gives
$|E_s[T_j]|\le2+2/(j^2-1)$. The panel error is therefore
\[
 h\sum_{\substack{j\ge2s\\j\ {\rm even}}}2\beta_j E_s[T_j].
\]
To estimate these higher-order coefficients, we sample an inner contour of parameter $1<\rho<\rho_2$
at $N$ roots of unity multiplied by $\rho$. The difference between the rescaled Fourier coefficient
of the samples and $\beta_j$ is
\[
 \sum_{\ell\ne0}\beta_{j+\ell N}\rho^{\ell N}+\rho^{-j}\widehat e_j ,
\]
where $e$ is the vector of errors and the Fourier transform is normalized by $1/N$. Geometric sums
on the outer ellipse bound the first term and the omitted tail. By Parseval's identity,
$\sum_{j=0}^{N-1}|\widehat e_j|^2\le N^{-1}\sum_i\delta_i^2$ if $|e_i|\le\delta_i$.
For the noise we retain the modes $j\le N/2-1$ and bound the rest of the tail analytically.
After multiplication by the Gaussian weights in the preceding expression, Cauchy--Schwarz bounds the
noise. For the 20-node panel we use 232 samples on the ellipse with $\rho=1.64$. For the six-node
panel, 24 independent samples represent a transform of length 46.

Under the stated sample tolerances, the bounds for the circle and the six- and twenty-point panels
are approximately $7.7\,10^{-10}$, $8.8\,10^{-11}$, and $1.7\,10^{-11}$, respectively.
The circle and short panel together give $8.56111216\,10^{-10}$. Adding the archived value for
the long panel, $1.69294096\,10^{-11}$, gives about $8.73040626\,10^{-10}$. The reported figure
of $8.7\,10^{-10}$ was rounded to the nearest value. Rounding upward gives $8.74\,10^{-10}$ [D if C].
This bound covers only the series and the integration rules under the stated tolerances.
To enclose $\cerr$, we must also propagate the errors at the 43 nodes of the calculation and in
the boundary terms. When a tolerance $\delta_i$ comes from grid refinement, its value is a
numerical estimate [N].
